% ECONOMICS_PROBLEM: {"id": "C73-20", "title": "Polynomial-time solution of simple stochastic games", "jel_code": "C73", "source_id": "OP-0163"}
% !TeX program = xelatex
\documentclass[11pt,a4paper]{article}
\usepackage[margin=23mm]{geometry}
\usepackage{fontspec}
\setmainfont{Latin Modern Roman}
\usepackage{amsmath,amssymb,mathtools}
\usepackage{xcolor,enumitem,booktabs,longtable,array,xurl,hyperref}
\definecolor{muted}{HTML}{53636C}
\hypersetup{colorlinks=true,urlcolor=blue,linkcolor=blue}
\setlength{\parindent}{0pt}
\setlength{\parskip}{5pt}
\newcommand{\R}{\mathbb R}
\newcommand{\E}{\mathbb E}
\newcommand{\N}{\mathbb N}
\newcommand{\ind}{\mathbf 1}
\newcommand{\Var}{\operatorname{Var}}
\newcommand{\Cov}{\operatorname{Cov}}
\newcommand{\supp}{\operatorname{supp}}
\DeclareMathOperator*{\argmin}{arg\,min}
\DeclareMathOperator*{\argmax}{arg\,max}
\newcommand{\entrymeta}[3]{{\sffamily\small\color{muted}\textbf{\texttt{#1}}\quad|\quad #2\par #3\par}}
\newcommand{\Problemref}[1]{\texttt{#1}}
\newcommand{\JELref}[2]{\texttt{#2} (JEL \texttt{#1})}
\begin{document}
\section*{Polynomial-time solution of simple stochastic games}
\textbf{Problem ID:} \texttt{C73-20}\quad \textbf{JEL:} \texttt{C73}\par
% BEGIN ECONOMICS STATEMENT
\label{problem:OP-0163}
% GAME_THEORY_PROBLEM: {"local_id": "GT-033", "original_number": 33, "kind": "P", "entry_kind": "statement_only", "published": false, "review_required": true, "category": "game_theory"}
\label{game:GT-033}
{\sffamily\small\color{muted}
{\sffamily\small\color{muted}\textbf{GT-033} | Stochastic games and computational complexity\par P | Source-reported polynomial-time question\par}
\par}
\subsection*{Definitions and mathematical statement}
A simple stochastic game has a finite directed graph with vertices controlled by Max, controlled by Min, or random, together with absorbing vertices $0$ and $1$. Every other vertex has two listed successors. At a controlled vertex its owner chooses the successor; at a random vertex each listed successor is chosen with probability $1/2$. Max's payoff is the indicator of ever reaching vertex $1$, and Min minimizes its expectation. Nonterminating paths have payoff zero. With arbitrary behavioral strategies, put
\[
 v(s)=\sup_{\sigma}\inf_{\tau}\Pr_{s,\sigma,\tau}(\text{reach }1).
\]
Given the explicitly represented graph, initial vertex $s$, and binary-encoded $q\in\mathbb Q\cap[0,1]$, the question is whether a deterministic algorithm decides $v(s)\geq q$ in time polynomial in the entire input bit length.
\subsection*{Short English statement}
Can the exact value-threshold problem for simple stochastic games be solved in polynomial time?
\subsection*{Variants and scope boundaries}
The input size includes the threshold encoding. A running time polynomial in $1/\varepsilon$ for approximate values is not polynomial-time exact solution. Stopping promises, broader rational transition kernels, and discounted variants must be related by explicit reductions before being called equivalent. An algorithm may output optimal pure stationary choices as a related search target, but exact rational value output must also account for its bit length.
\subsection*{Source relationship and status}
[S1]'s introduction retains simple stochastic game polynomial-time solvability as open and discusses its relationship with zero-sum discounted perfect-information games. The general-sum stationary-equilibrium classification in that paper does not resolve this zero-sum problem.
\subsection*{References}
\begingroup\footnotesize\raggedright
\textbf{[S1]} Kristoffer Arnsfelt Hansen and Xinhao Nie, \emph{On the Complexity of Stationary Nash Equilibria in Discounted Perfect Information Stochastic Games}, arXiv:2510.11550v1 (2025). Locator: Section 1, paragraphs on the polynomial equivalence with simple stochastic games and their unresolved complexity. \url{https://arxiv.org/html/2510.11550v1}
\par\endgroup

\subsection*{Common conventions: Source mathematical conventions}

\textbf{Finite games.} For a finite set $A$, $\Delta(A)$ is its probability
simplex. Players choose independent mixed strategies in Nash equilibrium;
correlation is allowed only when explicitly part of the solution concept.
In a game with payoff functions $u_i$ and strategy profile $x$, an additive
$\varepsilon$ Nash equilibrium satisfies
\[
 u_i(x)\geq u_i(x_i',x_{-i})-\varepsilon
 \quad\text{for every player $i$ and every admissible deviation $x_i'$}.
\]
For numerical approximation, payoffs are normalized as locally specified.
An $\varepsilon$ payoff residual is distinct from distance at most
$\varepsilon$ to the set of exact equilibria. Point convergence, convergence
of empirical play, and convergence in distance to an equilibrium set are
also distinct.

\textbf{Computational representations.} Unless stated otherwise, a complexity
question takes rational data with binary encoding; polynomial time is measured
in total bit length. A fixed-accuracy polynomial algorithm is not necessarily
polynomial in $1/\varepsilon$, and neither is necessarily polynomial in
$\log(1/\varepsilon)$. Succinct games and value, demand, or sampling oracles
must declare their interfaces. Exact real solutions require an explicit
representation; an unspecified list of arbitrary real numbers is not a
finite computational output. A claimed algorithmic barrier is meaningful
only for its specified model and reductions.

\textbf{Dynamic games.} Histories, information, observation, and allowable
randomization are part of the model. A stationary strategy depends only on
the current state; a positional strategy in a deterministic graph game
chooses one action at each controlled vertex. Nash equilibrium at an initial
state and equilibrium after every history are different requirements.
An expectation of a pathwise $\liminf$ payoff, a $\liminf$ of expected averages,
a limiting expected average, and a uniform finite-horizon equilibrium are
different criteria. None is silently substituted for another.

\textbf{Allocation and mechanisms.} A complete allocation assigns every item
exactly once unless otherwise stated. Zero-valued goods, negative values,
capacity constraints, feasibility, lotteries, and copies of goods affect
fairness definitions and are specified locally. Dominant-strategy incentive
compatibility, Bayesian incentive compatibility, universal truthfulness,
and truthfulness in expectation impose different inequalities. Whenever
an objective ratio has zero denominator, its convention is stated or those
instances are excluded.

\textbf{Infinite-dimensional models.} Expectations, integrals, stochastic
equations, and optimization objectives require their local regularity,
integrability, filtrations, and admissible domains. A backward--forward
mean-field system is a boundary-value problem; it is not an initial-value
semiflow without an additional construction. Long-time, large-population,
and vanishing-noise limits require an order of limits and a convergence
topology. If a minimum need not be attained, the objective is its infimum
or supremum and the approximation requirement is stated separately.
% END ECONOMICS STATEMENT
\end{document}
